% Consolidated textbook-ch04 from a4681eacb57a1f9fed79c1da26e5c6d0418a7d1c; source fragments preserved.
% 2026-11-24 八批严格复核第2批：本文件旧片段哈希对应历史版；当前内容含局部二校修订。

% === BEGIN textbook-ch04-choices.tex blob d6381d0b3c4bd07ba577f142939827cce7c0ee35 ===
% 《大学物理简明教程》（第5版）第4章选择题 4.1.1—4.1.12。
% 书内页/PDF页100—102；依据扫描题干和图形，非官方改述+原创推导。
% 自检不等于独立学术审定。
\subsection{选择题：简谐运动、波及干涉}

\sourceproblem{4.1.1}{100}{100}
\problemtext 一质点做简谐运动，$t=0$时$x=-A/2$且向$x$正方向运动；从旋转矢量法判断初相位。
\workings
以$x=A\cos(\omega t+\varphi)$表示，初始条件给
$\cos\varphi=-1/2$，而
$v(0)=-A\omega\sin\varphi>0$，所以$\sin\varphi<0$。
\[
\boxed{\varphi=\frac{4\pi}{3}\pmod{2\pi}}
\]
旋转矢量位于第三象限（与图示对应的选项为B）。
\crosscheck 若误取$2\pi/3$，速度会指向负$x$，不满足题设方向。

\sourceproblem{4.1.2}{100}{100}
\problemtext 两同周期简谐振动的位移—时间曲线如原图4.1.2，曲线(1)初始位于负最大位移，而曲线(2)先于(1)达到正最大位移。判断(1)相对于(2)的相位。
\workings
观察同一正峰位置，曲线(1)向右滞后四分之一周期$\Delta t=T/4$，故
\[
\boxed{\varphi_1-\varphi_2=-\omega\Delta t=-\frac{\pi}{2}}
\]
（模$2\pi$），即(1)落后(2) $\pi/2$，原题选A。
\crosscheck 对同频曲线使用峰间\emph{水平时间差}，不是直接比较不同峰的位移高低。

\sourceproblem{4.1.3}{100--101}{100--101}
\problemtext 做简谐运动的质点，周期为$T$；在$x=A/2$处向$x$正方向运动，问此后第一次到达$x=A$需要多长时间。
\workings
由$\cos\varphi=1/2$且$v=-A\omega\sin\varphi>0$，
取$\varphi=-\pi/3$。到达正极端的相位为$0$，变化$\Delta\theta=\pi/3$，故
\[
\boxed{\Delta t=\frac{\pi/3}{2\pi/T}=\frac T6}
\]
（原题选C）。
\crosscheck 非线性匀速位移法$A/2v$不适用，必须以相位差计算。

\sourceproblem{4.1.4}{101}{101}
\problemtext 连接同一轻弹簧的物体置于竖直方向或光滑斜面上，在弹簧始终服从胡克定律且物体受适当约束时，能否做简谐振动？
\workings
沿允许运动的坐标$s$，重力分力为常数$F_g$，弹力为$-ks$，
\[
m\ddot s=F_g-ks.
\]
设平衡位置$s_e=F_g/k$，令$\xi=s-s_e$，便有
$m\ddot\xi+k\xi=0$，即$\xi$为简谐运动。竖直和光滑斜面两情形均成立（选C）。
\crosscheck 这以弹簧保持连接、无松弛、无碰撞和线性胡克定律为前提；重力只是移动平衡位置，不改变$\omega=\sqrt{k/m}$。

\sourceproblem{4.1.5}{101}{101}
\problemtext $x=A\cos(\omega t+\varphi)$，在$t=T/2$时瞬时速度是多少？
\workings
$\omega T=2\pi$，
\[
v(T/2)=-A\omega\sin(\pi+\varphi)
=\boxed{A\omega\sin\varphi}.
\]
原题选B。
\crosscheck 半周期后速度与初速反号；初速$-A\omega\sin\varphi$。

\sourceproblem{4.1.6}{101}{101}
\problemtext 质点以频率$\nu$作简谐振动，问动能随时间的变化频率。
\workings
由$v=-A\omega\sin(\omega t+\varphi)$，
\[
K=\frac12mA^2\omega^2\sin^2(\omega t+\varphi)
=\frac14mA^2\omega^2[1-\cos(2\omega t+2\varphi)].
\]
因此能量变化角频率是$2\omega$，频率
$\boxed{\nu_K=2\nu}$，原题选C。
\crosscheck 质点经过一个周期有两次速度最大、两次速度为零。

\sourceproblem{4.1.7}{101}{101}
\problemtext 在简谐振动中，动能与势能相等的位置是什么？
\workings
总能量$E=\frac12kA^2$，势能$U=\frac12kx^2$，
$K=E-U$。令$K=U$，得
\[
kx^2=\frac12kA^2,\qquad
\boxed{x=\pm\frac{A}{\sqrt2}}.
\]
原题选D。
\crosscheck 两位置关于平衡点对称，能量均为总能量一半。

\sourceproblem{4.1.8}{101}{101}
\problemtext 弹簧振子质量、弹簧劲度系数不变，振幅变为原来的两倍。比较周期、能量、最大速率与最大加速度。
\workings
\[
T=2\pi\sqrt{\frac mk},\quad
E=\frac12kA^2,\quad
v_{\max}=\omega A,\quad
a_{\max}=\omega^2 A.
\]
令$A\mapsto2A$后，$T$不变，$E$变为$4E$，
$v_{\max}$及$a_{\max}$均变为2倍。原题选C。
\crosscheck 线性弹簧振子中周期与振幅无关，不适用于任意大幅非线性振动。

\sourceproblem{4.1.9}{101}{101}
\problemtext 波沿$x$正方向以速率$u$传播；距$O$点$l$的$P$点满足$y_P=A\cos\omega t$，求$O$点和整条波的波函数。
\workings
$P$在$x=l$，正向行波一般写为$y(x,t)=A\cos[\omega t-k(x-l)]$，$k=\omega/u$。
因此
\[
\boxed{y(x,t)=A\cos\omega\left(t+\frac lu-\frac xu\right)},\qquad
\boxed{y_O(t)=A\cos\omega\left(t+\frac lu\right)}.
\]
原题对应选C。
\crosscheck 代入$x=l$恢复$P$点已知振动；固定相位条件$x-ut=\mathrm{const}$为正向传播。

\sourceproblem{4.1.10}{101}{101}
\problemtext 一机械波$y=0.05\cos(6\pi t+0.06\pi x)$（SI）。判断波长、波速、周期与传播方向。
\workings
由$\omega=6\pi\,\mathrm{rad/s}$、$k=0.06\pi\,\mathrm{rad/m}$可得
\[
f=3\,\mathrm{Hz},\quad
\boxed{T=\frac13\,\mathrm s},\quad
\lambda=\frac{2\pi}{k}=\frac{100}{3}\,\mathrm m,\quad
u=\frac{\omega}{k}=100\,\mathrm{m/s}.
\]
相位$\omega t+kx$为常量时$x$随$t$减小，故沿$x$负方向传播。
四个选项中仅周期为$1/3\,\mathrm s$正确（C）。
\crosscheck 不能将方程中$0.06\pi$直接当成$2\pi/\lambda$之外的“波长”。

\sourceproblem{4.1.11}{101--102}{101--102}
\problemtext 两相干波源$S_1,S_2$相位分别为$\varphi_1,\varphi_2$，到$P$距离$r_1,r_2$，波长均为$\lambda$。写出在$P$干涉加强的相位条件。
\workings
两波到达$P$的相位分别为
$\theta_1=\omega t+\varphi_1-2\pi r_1/\lambda$与
$\theta_2=\omega t+\varphi_2-2\pi r_2/\lambda$。
相位差
\[
\Delta\theta=\theta_2-\theta_1
=\varphi_2-\varphi_1+\frac{2\pi(r_1-r_2)}{\lambda}.
\]
因此干涉加强的条件是
\[
\boxed{\varphi_2-\varphi_1+\frac{2\pi(r_1-r_2)}{\lambda}=2n\pi},
\quad n\in\mathbb Z.
\]
原题选D。
\crosscheck $r_1=r_2$且$\varphi_1=\varphi_2$时必加强；若两源反相需半整数波长差。

\sourceproblem{4.1.12}{102}{102}
\problemtext 弹性介质中传播的平面简谐行波，一质元从平衡位置移动到最大位移位置的过程中，它的机械能如何变化？
\workings
对理想线性无耗散行波，局部动能密度和弹性势能密度同相变化。以
$y=A\cos(\omega t-kx)$为例，
\[
\mathcal E_k=\frac12\rho(\partial_t y)^2
=\frac12\rho A^2\omega^2\sin^2(\omega t-kx),
\]
对一维弦，以上$\rho$须取线密度（单位$\mathrm{kg/m}$），$\mathcal E_k$为单位长度动能；张力记为$F_T$，单位长度弹性势能为
$\mathcal E_p=\frac12F_T(\partial_x y)^2
=\frac12\rho A^2\omega^2\sin^2(\omega t-kx)$，其中$F_T/\rho=u^2$。若采用三维介质的体积能量密度，则$\rho$为体密度，弹性系数须改为相应的弹性模量（单位Pa），不可直接代入弦张力。
质元由平衡位移达到最大位移时，两个局部能量密度均由最大变为零，
能量通过介质传给相邻部分。原题选D。
\crosscheck 传播波的一小段介质并非孤立弹簧振子，局部总能量不必守恒；静止正弦驻波的能量分布则不同。
% === END textbook-ch04-choices.tex ===

% === BEGIN textbook-ch04-blanks.tex blob 9af0d19690440637a8d0e1111c9e32f5177aea19 ===
% 《大学物理简明教程》第5版，第4章填空题4.2.1—4.2.13
% 书内/PDF第102—103页；源题图以扫描原书为准。
\subsection{填空题：简谐振动与波}

\sourceproblem{4.2.1}{102}{102}
\problemtext 振幅$A=4\,\mathrm{cm}$、周期$T=2\,\mathrm s$，在$t=0$第一次经过$x=-2\,\mathrm{cm}$且沿负方向运动。求此后第二次经过该位置的时刻。
\workings
令$x=A\cos(\omega t+\varphi)$，$\omega=2\pi/T=\pi\,\mathrm{rad/s}$。
初相位满足$\cos\varphi=-1/2$且$\sin\varphi>0$，取$\varphi=2\pi/3$。
下次经过$x=-A/2$时$\omega t+\varphi=4\pi/3$，故
\[
\boxed{t=\frac{(4\pi/3)-(2\pi/3)}{\pi}=\frac23\,\mathrm s}.
\]
\crosscheck 该次速度沿正方向；同方向再次经过则需一个完整周期$2\,\mathrm s$。

\sourceproblem{4.2.2}{102}{102}
\problemtext 原题4.2.2的简谐振动$x$—$t$图标有$a,b,c,d,e,f$各点。判别位移零、速度$+\omega A$、加速度零、弹力零的点，以及速度零、加速度$-\omega^2A$、弹力$-kA$的点。
\workings
由
$v=\dot x,\ a=-\omega^2x,\ F=-kx$，
在$x=0$且曲线向上穿过横轴的点$d$，有$v=+\omega A,a=F=0$。
在$x=+A$的正峰点$a$和$e$，有$v=0,a=-\omega^2A,F=-kA$。
\[
\boxed{\text{第一组：}d;\qquad\text{第二组：}a\text{或}e}.
\]
\crosscheck 正峰出现于每周期一次，若题面空格只接受一个字母，可填$a$（或下一周期的$e$）；峰值状态不是唯一时刻。

\sourceproblem{4.2.3}{102}{102}
\problemtext 质点周期$T$、振幅$A$沿$x$轴简谐运动。(a) $t=0$位于原点且沿正方向；(b) $t=0$位于$x=A/2$且沿负方向。分别求振动方程。
\workings
记$\omega=2\pi/T$。
(a) $x(0)=0,\ v(0)>0$，可写为
\[
\boxed{x_a=A\sin(2\pi t/T)}.
\]
(b) $\cos\varphi=1/2,\ v(0)=-A\omega\sin\varphi<0$，故$\varphi=\pi/3$，
\[
\boxed{x_b=A\cos(2\pi t/T+\pi/3)}.
\]
\crosscheck 两式的初始速度符号分别为$+A\omega$与$-\sqrt3 A\omega/2$。

\sourceproblem{4.2.4}{102}{102}
\problemtext 弹簧劲度系数$k=0.32\,\mathrm{N/m}$，重物质量$m=0.02\,\mathrm{kg}$。求固有频率与周期。
\workings
\[
\omega_0=\sqrt{k/m}=\sqrt{16}=4\,\mathrm{rad/s},
\qquad
\boxed{\nu_0=\frac{\omega_0}{2\pi}=\frac2\pi\,\mathrm{Hz}},
\quad
\boxed{T=\frac{2\pi}{\omega_0}=\frac\pi2\,\mathrm s}.
\]
\crosscheck 周期约$1.571\,\mathrm s$，频率约$0.637\,\mathrm{Hz}$，二者互为倒数。

\sourceproblem{4.2.5}{102}{102}
\problemtext 原图4.2.5中，两个简谐振动$x_1,x_2$的初始位移均为零且向正方向运动；比较频率$\nu_1:\nu_2$、最大加速度$a_1:a_2$以及初始速率$v_{10}:v_{20}$。
\workings
按图中较密集的$x_1$曲线与较疏的$x_2$曲线，前者周期约为后者的一半，振幅绘制为相同大小。若将该图按$T_1=T_2/2$、$A_1=A_2$解读，则
\[
\boxed{\nu_1:\nu_2=2:1},\quad
a_{\max,1}:a_{\max,2}
=A_1\omega_1^2:A_2\omega_2^2=\boxed{4:1},
\]
\[
\boxed{v_{10}:v_{20}=A_1\omega_1:A_2\omega_2=2:1}.
\]
\crosscheck \textbf{图像比例仍待高分辨率二校。}本结果依赖振幅相同、两者初速度取正且图示周期比恰为$1:2$；若原图仅为示意并未给出严格尺度，不能把上述比值当作独立测量结果。

\sourceproblem{4.2.6}{102}{102}
\problemtext 原图4.2.6的位移曲线单位为cm与s；振幅为$2$ cm，初始位移为$-1$ cm且向负方向，$t=1\,\mathrm s$为正极值，$t=0.25\,\mathrm s$为负极值。求振动方程。
\workings
相邻正负极值间隔$0.75\,\mathrm s=T/2$，所以
$T=1.5\,\mathrm s$、$\omega=4\pi/3\,\mathrm{rad/s}$。
$\cos\varphi=-1/2$且初速为负，$\varphi=2\pi/3$，得
\[
\boxed{x(t)=2\cos\left(\frac{4\pi}{3}t+\frac{2\pi}{3}\right)\,\mathrm{cm}}.
\]
\crosscheck 代入$t=0.25\,\mathrm s$得$x=-2$ cm，$t=1\,\mathrm s$得$x=+2$ cm，符合图示极值。

\sourceproblem{4.2.7}{102}{102}
\problemtext 简谐波频率$\nu=5\times10^4\,\mathrm{Hz}$、波速$u=1.5\times10^3\,\mathrm{m/s}$，沿传播路径相距$5\times10^{-3}\,\mathrm m$两点的相位差。
\workings
$\lambda=u/\nu=0.03\,\mathrm m$，同一时刻沿传播方向两点相位差的大小
\[
\boxed{|\Delta\varphi|=\frac{2\pi\Delta x}{\lambda}
=\frac{2\pi(0.005)}{0.03}=\frac\pi3}.
\]
\crosscheck 指明相位有向差时需区分两点先后位置，题只问相位差大小。

\sourceproblem{4.2.8}{102}{102}
\problemtext 平面简谐波频率$100\,\mathrm{Hz}$、传播速率$300\,\mathrm{m/s}$，两点同一时刻的相位差大小为$\pi/3$；问相距多少。
\workings
波长$\lambda=u/\nu=3\,\mathrm m$；在沿传播方向且取最近的两点这一条件下：
\[
|\Delta x|_{\min}=\frac{\pi/3}{2\pi}\lambda
=\boxed{0.50\,\mathrm m}.
\]
\crosscheck 若相位差理解为模$2\pi$，还存在
$\Delta x=3n\pm0.5\,\mathrm m$的其他非负距离；若题面不强调“最近”，单一距离并非数学上唯一。在三维平面波中，相位差约束的是两点间距沿传播方向的投影，非任意两点的欧氏距离。本稿以最小正距离填写并保留此限定。

\sourceproblem{4.2.9}{102--103}{102--103}
\problemtext 一横波的波函数为
$y=0.02\sin[2\pi(100t-0.4x)]$（SI），求振幅、波长、频率和传播速度。
\workings
对照$y=A\sin[2\pi(\nu t-x/\lambda)]$，
\[
\boxed{A=0.02\,\mathrm m},\quad
\boxed{\lambda=1/0.4=2.5\,\mathrm m},\quad
\boxed{\nu=100\,\mathrm{Hz}},\quad
\boxed{u=\lambda\nu=250\,\mathrm{m/s}}.
\]
\crosscheck 相位$\nu t-x/\lambda$为常数时$x$随$t$增大，表明沿正$x$传播。

\sourceproblem{4.2.10}{103}{103}
\problemtext 机械波产生的必要条件是什么？
\workings
需要能提供扰动的\textbf{波源（振源）}（单次扰动也能产生脉冲波，周期性仅是持续周期波的条件），以及能够把振动通过相互作用逐点传递的\textbf{弹性介质}。
\[
\boxed{\text{波源（振源）与传播介质}}
\]
\crosscheck 声波在真空中不能传播；“有振源”而没有能相互耦合的介质，不会形成经典机械波。

\sourceproblem{4.2.11}{103}{103}
\problemtext 两列机械波叠加产生稳定干涉现象的基本条件。
\workings
稳定干涉要求
\[
\boxed{\text{频率相同（或固定频差为零）、相位差恒定、振动方向相同或有公共分量}}.
\]
\crosscheck 对纯横向正交偏振且无公共振动分量的两波，在单一方向的能量观测中无常规明暗干涉项；相位随机游走时干涉条纹不稳定。

\sourceproblem{4.2.12}{103}{103}
\problemtext 波沿$x$轴正向传播，$u=0.08\,\mathrm{m/s}$，$\lambda=0.20\,\mathrm m$。$t=0$的波形如图4.2.12：原点位移为零且沿$x$增大时位移先为正，振幅为$0.04\,\mathrm m$。求波函数。
\workings
由波速与波长求得波数、角频率：
\[
k=\frac{2\pi}{\lambda}=10\pi\,\mathrm{rad/m},
\qquad
\omega=uk=0.8\pi\,\mathrm{rad/s}.
\]
$t=0$时空间函数为
\[
y(x,0)=0.04\sin(10\pi x).
\]
沿正$x$传播，将初始函数中的$x$替换成$x-ut$：
\[
\boxed{y(x,t)=0.04\sin(10\pi x-0.8\pi t)\ \mathrm m}.
\]
\crosscheck $y(0,0)=0$且$\partial_x y(0,0)>0$；
固定波峰位置满足$x-ut=\mathrm{const}$，故其速度为$+0.08\,\mathrm{m/s}$。

\sourceproblem{4.2.13}{103}{103}
\problemtext 平面简谐波经截面积$S=3.00\times10^{-2}\,\mathrm{m^2}$的空气管传播，$10$ s通过该截面的波能为$2.70\times10^{-2}$ J。求平均能流密度。
\workings
平均功率$\bar P=E/\Delta t=2.7\times10^{-3}\,\mathrm W$，
除以面积得
\[
\boxed{\bar I=\frac{E}{S\Delta t}
=\frac{2.70\times10^{-2}}{(3.00\times10^{-2})(10)}
=0.090\,\mathrm{W/m^2}}.
\]
\crosscheck 单位$\mathrm J/(\mathrm{m^2\,s})=\mathrm{W/m^2}$。
% === END textbook-ch04-blanks.tex ===

% === BEGIN textbook-ch04-solutions-001-010.tex blob 8fc78e56fcbd1af3f185fa4085ebeaeb8991778d ===
% 《大学物理简明教程》第5版，第4章解答题4.3.1—4.3.10。
% 原书印刷页与扫描PDF第103—104页；推导单独建立而非复制原书答案。
\subsection{解答题：简谐运动与振动合成}

\sourceproblem{4.3.1}{103}{103}
\problemtext 小球位于竖直平面内的光滑圆弧形轨道最低点附近，轨道半径$R$；作小幅运动时证明它为简谐振动，并求周期。
\workings
令偏离竖直向下的圆心角为$\theta$。轨道切向重力分力为$-mg\sin\theta$，故
\[
mR\ddot\theta=-mg\sin\theta.
\]
此处小球视为质点滑动而非有限尺寸刚体纯滚动。小角度$|\theta|\ll1$时$\sin\theta\simeq\theta$，
\[
\ddot\theta+\frac gR\theta=0,\qquad
\boxed{\omega_0=\sqrt{g/R}},\quad
\boxed{T=2\pi\sqrt{R/g}}.
\]
\crosscheck 小球若沿轨道大幅运动，$\sin\theta$的非线性会使周期依赖振幅，以上只成立于小幅近似。

\sourceproblem{4.3.2}{103}{103}
\problemtext 一弹簧振子质量$m=10\times10^{-3}\,\mathrm{kg}$，位移方程
$x=0.1\cos(8\pi t+2\pi/3)$（SI）。求周期、振幅、初相、速度及加速度最大值、最大回复力、总能量、时间平均动能与势能、动势能相等位置，以及$t_1=1\,\mathrm s$和$t_2=5\,\mathrm s$相位差。
\workings
$A=0.1\,\mathrm m,\ \omega=8\pi\,\mathrm{rad/s},\
\varphi_0=2\pi/3$，故
\[
\boxed{T=\frac14\,\mathrm s},\quad
\boxed{v_{\max}=A\omega=0.8\pi\,\mathrm{m/s}},\quad
\boxed{a_{\max}=A\omega^2=6.4\pi^2\,\mathrm{m/s^2}}.
\]
$k=m\omega^2=0.64\pi^2\,\mathrm{N/m}$，
\[
\boxed{F_{\max}=kA=0.064\pi^2\,\mathrm N},\quad
\boxed{E=\frac12kA^2=0.0032\pi^2\,\mathrm J}.
\]
时间平均满足$\langle K\rangle=\langle U\rangle=E/2=0.0016\pi^2\,\mathrm J$；
$K=U$时$x=\pm A/\sqrt2=\pm0.1/\sqrt2\,\mathrm m$。
两时刻的相位差
\[
\boxed{\Delta\varphi=\omega(t_2-t_1)=32\pi
\equiv0\pmod{2\pi}}.
\]
\crosscheck 平均能量取整周期或长时间平均，势能零点取平衡位置。经过4秒恰好16个完整振动周期，初末状态相同；$\varphi$未约化时的“累计相位差”为$32\pi$。

\sourceproblem{4.3.3}{103}{103}
\problemtext 沿$x$轴的简谐振动振幅$A$、周期$T$，采用余弦式$x=A\cos(\omega t+\varphi)$。分别求四种$t=0$条件的初相：(1) $x_0=-A$；(2)过平衡点向正向运动；(3)经过$x=A/2$向负向；(4)经过$x=-A/\sqrt2$向正向。
\workings
$\omega=2\pi/T$，速度$v=-A\omega\sin(\omega t+\varphi)$，由位移确定余弦、由速度符号确定正弦。
\[
\begin{array}{c|c|l}
\text{条件}&\varphi\pmod{2\pi}&x(t)\\ \hline
(1)&\pi&A\cos(\omega t+\pi)\\
(2)&-\pi/2&A\cos(\omega t-\pi/2)\\
(3)&\pi/3&A\cos(\omega t+\pi/3)\\
(4)&-3\pi/4&A\cos(\omega t-3\pi/4)
\end{array}
\]
\crosscheck (3) $\sin(\pi/3)>0$，速度负；(4) $\cos(-3\pi/4)=-1/\sqrt2$，$\sin(-3\pi/4)<0$，速度正。初相可以加任意整数倍$2\pi$。

\sourceproblem{4.3.4}{103}{103}
\problemtext 一振子质量$0.01\,\mathrm{kg}$，振幅$A=24\,\mathrm{cm}$，周期$T=4.0\,\mathrm s$。$t=0$位于$x=+A$。
(1) 求$t=0.5\,\mathrm s$的位置和受力；(2) 最短多久从初始位置到达$x=12\,\mathrm{cm}$；(3) 从$x=+12\,\mathrm{cm}$第一次到达$x=-12\,\mathrm{cm}$的最短时间。
\workings
$\omega=2\pi/T=\pi/2\,\mathrm{rad/s}$，
$x(t)=0.24\cos(\pi t/2)\,\mathrm m$。
(1) $t=0.5$：
\[
\boxed{x=0.24\cos(\pi/4)=0.12\sqrt2\,\mathrm m\approx0.1697\,\mathrm m},
\]
\[
\boxed{F=-m\omega^2x
=-0.01\left(\frac\pi2\right)^2(0.12\sqrt2)
\approx-4.19\times10^{-3}\,\mathrm N}.
\]
(2) 初始在正极端，$\cos(\omega t)=1/2$，首次时间
\[
\boxed{t_{A\to A/2}=\frac{\pi/3}{\pi/2}=\frac23\,\mathrm s}.
\]
(3) 沿继续负向运动从$+A/2$到$-A/2$，相位由$\pi/3$到$2\pi/3$，
\[
\boxed{\Delta t=\frac{(\pi/3)}{\pi/2}=\frac23\,\mathrm s}.
\]
\crosscheck 两问时间相等是此特定相位几何的结果，不是任意振动区段的普遍性质。

\sourceproblem{4.3.5}{103}{103}
\problemtext 光滑水平桌面上，一质量$M$的木块连接劲度系数$k$的轻弹簧且初始静止、弹簧为自然长。质量$m$的子弹以水平速率$u_0$射入木块并滞留，求碰后的振动方程及在位移$x=A/2$处的动能、势能。
\workings
短时黏连碰撞中平动动量守恒，公共速率
$V=mu_0/(M+m)$。随后无摩擦弹簧振动：
\[
\omega=\sqrt{\frac{k}{M+m}},\qquad
A=\frac{V}{\omega}
=\frac{mu_0}{\sqrt{k(M+m)}}.
\]
以碰撞瞬间为$t=0$、子弹入射方向为$x$正向，
\[
\boxed{x(t)=A\sin(\omega t)}.
\]
总机械能$E=\frac12kA^2$，当$x=A/2$，
\[
\boxed{U=\frac18kA^2},\qquad
\boxed{K=E-U=\frac38kA^2}.
\]
\crosscheck \emph{不能}将子弹碰前动能直接等同于碰后振动总能量；完全非弹性碰撞已损失机械能。

\sourceproblem{4.3.6}{104}{104}
\problemtext 竖直轻弹簧挂质量$1.0\,\mathrm g$的小物体时伸长$4.9\,\mathrm{cm}$。改由弹簧与质量$8.0\,\mathrm g$的小球构成振子：从平衡位置向下拉$1.0\,\mathrm{cm}$，在$t=0$赋予向上的初速度$5.0\,\mathrm{cm/s}$。求周期和振动方程。
\workings
静伸长求弹簧劲度：
\[
k=\frac{(0.001)(9.8)}{0.049}=0.20\,\mathrm{N/m}.
\]
振子$m=0.008\,\mathrm{kg}$，角频率与周期
\[
\omega=\sqrt{k/m}=5\,\mathrm{rad/s},\quad
\boxed{T=\frac{2\pi}{5}\,\mathrm s}.
\]
取向下为$x$正方向，则$x_0=+0.010\,\mathrm m,\ v_0=-0.050\,\mathrm{m/s}$。
\[
A=\sqrt{x_0^2+(v_0/\omega)^2}
=0.010\sqrt2\,\mathrm m,\quad
\cos\varphi=\frac1{\sqrt2},\ \sin\varphi=\frac1{\sqrt2}.
\]
所以
\[
\boxed{x(t)=0.010\sqrt2\cos(5t+\pi/4)\,\mathrm m}.
\]
\crosscheck $x(0)=+1$ cm、$\dot x(0)=-5$ cm/s；由另一物体的静伸长估算$k$时不能误把$1$ g当成本次运动质量。

\sourceproblem{4.3.7}{104}{104}
\problemtext 根据图4.3.7的两条$x$—$t$曲线，分别写出简谐运动方程。图(a)在$t=0$穿过$x=0$向上，$t=1\,\mathrm s$向下穿越零、$t=2\,\mathrm s$向上穿越零，振幅$10$ cm；图(b)在$t=0$位于$+5$ cm且向上，振幅$10$ cm，图示周期约为$2$ s。
\workings
(a) 正向过零到下一次同向过零为一个周期$2\,\mathrm s$，故
\[
\boxed{x_a(t)=10\sin(\pi t)\,\mathrm{cm}}.
\]
(b) 由$t=0$处$x/A=1/2$且曲线向上，余弦初相为$-\pi/3$；
按图周期$T=2$ s，
\[
\boxed{x_b(t)=10\cos(\pi t-\pi/3)\,\mathrm{cm}}.
\]
\crosscheck $x_b(0)=5$ cm且$\dot x_b(0)=+5\sqrt3\pi$ cm/s；
(b) 峰值发生于$t=1/3$ s。原图横轴刻度分辨率仍需单独图形二校。

\sourceproblem{4.3.8}{104}{104}
\problemtext 两个同方向、同频率简谐振动合振幅为$A=0.20\,\mathrm m$，合振动相位相对第一个振动\emph{超前}$\pi/6$，已知第一个振幅$A_1=0.173\,\mathrm m$。求第二个振幅以及两振动间的相位差。
\workings
以第一个振动的旋转矢量沿实轴，合矢量为$A e^{i\delta}$，$\delta=\pi/6$，
第二振动矢量为$A_2e^{i\Delta}=Ae^{i\delta}-A_1$。分别比较实虚部，
\[
A_2^2=(A\cos\delta-A_1)^2+(A\sin\delta)^2
=A^2+A_1^2-2AA_1\cos\delta,
\]
\[
\boxed{A_2=\sqrt{A^2+A_1^2-2AA_1\cos(\pi/6)}
\approx0.100\,\mathrm m}.
\]
从分量求相位差
\[
\boxed{\Delta=\operatorname{atan2}
\left(A\sin\frac\pi6,\ A\cos\frac\pi6-A_1\right)
\approx\frac{\pi}{2}}.
\]
\crosscheck $0.173\,\mathrm m$是约$\sqrt3(0.1)$ m，故若$A=0.2$ m且合矢量领先$\pi/6$，另一个振幅约0.1 m且与第一振动相差约$90^\circ$；按旋转矢量分解，而非误将给定的$\pi/6$当作两原振动相差。

\sourceproblem{4.3.9}{104}{104}
\problemtext 分别求两组同频同向振动合成后的振幅：
(1) $x_1=5\cos(3t+\pi/3)$ cm，$x_2=5\cos(3t+7\pi/3)$ cm；
(2) $x_1=5\cos(3t+\pi/3)$ cm，$x_2=5\cos(3t+4\pi/3)$ cm。
\workings
(1) 相位差为$2\pi$，完全同相，故$\boxed{A=10\,\mathrm{cm}}$。
(2) 相位差为$\pi$，幅度相同而相反，故$\boxed{A=0}$。
\crosscheck 两振动的角频率均为$3\,\mathrm{rad/s}$；只有同频情况下能用固定相量法。

\sourceproblem{4.3.10}{104}{104}
\problemtext 同一直线同频振动$x_1=0.4\cos(2t+\pi/6)$ m与
$x_2=0.3\cos(2t-5\pi/6)$ m，求合振幅、合初相和合振动方程。
\workings
两初相之差为$-\pi$，故两个相量反向；
振幅相减得$A=0.4-0.3=0.1$ m，合相位与第一个一致：
\[
\boxed{A=0.1\,\mathrm m},\qquad
\boxed{\varphi=\pi/6\pmod{2\pi}},\qquad
\boxed{x(t)=0.1\cos(2t+\pi/6)\,\mathrm m}.
\]
\crosscheck 在$t=0$时$x_1+x_2=(0.4-0.3)\cos(\pi/6)$，与合式完全一致。
% === END textbook-ch04-solutions-001-010.tex ===

% === BEGIN textbook-ch04-solutions-011-017.tex blob 4ab3ffb3c7b6e46d3700ce1d8f8c4f39b6bc4105 ===
% 《大学物理简明教程》第5版，第4章解答题4.3.11—4.3.17。
% 来源印刷/PDF第104—105页，波形图模型按原图定位；待学术二校。
\subsection{解答题：平面行波与相位}

\sourceproblem{4.3.11}{104}{104}
\problemtext 已知沿$x$轴传播的平面简谐波
$y=A\cos(Bt-Cx)$，$A,B,C>0$。求振幅、波速、频率、周期和波长；
原点的振动方程；以及沿传播方向相距$\lambda/2$两点的相位差。
\workings
与标准式$y=A\cos(\omega t-kx)$比较，$\omega=B,k=C$，故
\[
\boxed{A=A,\quad u=\frac BC,\quad
\nu=\frac B{2\pi},\quad T=\frac{2\pi}{B},\quad
\lambda=\frac{2\pi}{C}}.
\]
$x=0$处振动为$\boxed{y_O(t)=A\cos Bt}$。
同一时刻，相距$\lambda/2$的相位差大小
\[
\boxed{|\Delta\varphi|=C\lambda/2=\pi}.
\]
\crosscheck 固定相位$Bt-Cx=\mathrm{const}$的轨迹$x=(B/C)t+\mathrm{const}$，方向为$x$正向。

\sourceproblem{4.3.12}{104}{104}
\problemtext 沿细绳传播的简谐波
$y=0.05\cos(10\pi t-4\pi x)$（SI）。求：
(1) 各质点最大速度和加速度；
(2) $x=0.2\,\mathrm m$处$t=1\,\mathrm s$时的相位与原点在哪个时刻相同；
(3) 该相位在$t=1.25\,\mathrm s$时传播到哪里。
\workings
$\omega=10\pi\,\mathrm{rad/s},k=4\pi\,\mathrm{rad/m}$，
$u=\omega/k=2.5\,\mathrm{m/s}$。
\[
\boxed{v_{\max}=A\omega=0.5\pi\,\mathrm{m/s}},\qquad
\boxed{a_{\max}=A\omega^2=5\pi^2\,\mathrm{m/s^2}}.
\]
$x=0.2,t=1$的连续相位$\theta=10\pi-0.8\pi=9.2\pi$ rad
（模$2\pi$为$1.2\pi$）。原点有相同相位时
$10\pi t_0=9.2\pi+2n\pi$，故
\[
\boxed{t_0=0.92+0.2n\ \mathrm s,\quad n\in\mathbb Z}.
\]
追踪该同一连续相位的相位面（这里的相位不是波峰）到$t=1.25$，
\[
\boxed{x=0.2+u(1.25-1)=0.825\,\mathrm m}.
\]
\crosscheck 传播距离为$2.5\times0.25=0.625$ m。相位相同的\emph{其他周期副本}的坐标相隔整波长$\lambda=0.5$ m。

\sourceproblem{4.3.13}{104--105}{104--105}
\problemtext 原题4.3.13的某时刻空间波形沿$x$轴：$O$处过零向上，$A$为第一个正峰，$B$为随后的向下过零点，$C$为再下一个向上过零点。分别讨论波沿$+x$和$-x$传播时该时刻四点各自振动的相位。
\workings
以$A$到$C$间图像确定各点的\emph{空间}相位：
$\theta_s(O)=0,\theta_s(A)=\pi/2,\theta_s(B)=\pi,\theta_s(C)=2\pi$，故空间函数可写$y(x,0)=A_0\sin kx$（起点作适当选择）。
若波沿$+x$传播，
$y=A_0\cos(\omega t-kx+\pi/2)$；
在该时刻，以余弦式定义的四点时间相位为
\[
\boxed{\varphi_+(O)=\frac\pi2,\quad
\varphi_+(A)=0,\quad
\varphi_+(B)=-\frac\pi2,\quad
\varphi_+(C)=-\frac{3\pi}{2}}.
\]
若波沿$-x$传播，
$y=A_0\cos(\omega t+kx-\pi/2)$，
相应
\[
\boxed{\varphi_-(O)=-\frac\pi2,\quad
\varphi_-(A)=0,\quad
\varphi_-(B)=\frac\pi2,\quad
\varphi_-(C)=\frac{3\pi}{2}}.
\]
\crosscheck 相位仅在模$2\pi$意义下确定；本解固定采用\emph{时间余弦相位}，避免把波形的空间正弦相位直接当成质点的时间相位。

\sourceproblem{4.3.14}{104--105}{104--105}
\problemtext 沿$x$轴正方向传播的平面简谐波，波速$u=5\,\mathrm{m/s}$，波长$\lambda=2\,\mathrm m$。原点的振动曲线在$t=0$过零并向正位移运动，振幅$0.1\,\mathrm m$。求波函数、$t=0$的波形及距原点$0.5\,\mathrm m$处的振动曲线。
\workings
$f=u/\lambda=2.5\,\mathrm{Hz}$，
$\omega=5\pi\,\mathrm{rad/s}$，$k=\pi\,\mathrm{rad/m}$。
原点初相以正弦表述为零且初速为正，
\[
\boxed{y(x,t)=0.1\sin(5\pi t-\pi x)\,\mathrm m}.
\]
$t=0$时$\boxed{y(x,0)=-0.1\sin(\pi x)\,\mathrm m}$；
$x=+0.5$ m点的振动方程
\[
\boxed{y(0.5,t)=0.1\sin(5\pi t-\pi/2)
=-0.1\cos(5\pi t)\,\mathrm m}.
\]
\crosscheck 右行波在$t=0$空间原点的斜率为负，和原点随时间正速度相反：
$\partial_ty=-u\,\partial_xy$。

\sourceproblem{4.3.15}{105}{105}
\problemtext 图4.3.15给出$t=0$与$t=0.5\,\mathrm s$的波形(a)(b)，波沿$x$正方向传播。振幅$0.1\,\mathrm m$，(a)中相邻正峰间距约$4\,\mathrm m$，(b)相对(a)右移$1\,\mathrm m$。求波函数和$x=1\,\mathrm m$的$P$点振动方程。
\workings
由图可得$\lambda=4\,\mathrm m$、$k=\pi/2\,\mathrm{rad/m}$，
移速$u=1/0.5=2\,\mathrm{m/s}$，$\omega=ku=\pi\,\mathrm{rad/s}$。
$t=0$曲线在$x=0$向上穿过零，故
\[
\boxed{y(x,t)=0.1\sin\left(\frac\pi2x-\pi t\right)\,\mathrm m}.
\]
代入$x=1$，
\[
\boxed{y_P(t)=0.1\cos(\pi t)\,\mathrm m}.
\]
\crosscheck 源图只画有限波段，位移$1$ m取为向右的\emph{最小正}图示平移；相差整波长的平移给出等效波形。严格地，若无额外条件，右移可为$(1+4n)\,\mathrm m$（$n=0,1,2,\ldots$），$u_n=(2+8n)\,\mathrm{m/s}$，$\omega_n=(1+4n)\pi\,\mathrm{rad/s}$，$y_n(x,t)=0.1\sin[\pi x/2-(1+4n)\pi t]$ m，$y_{P,n}=0.1\cos[(1+4n)\pi t]$ m。上面的单一答案是最小正平移分支；需核原图有无额外约束。

\sourceproblem{4.3.16}{105}{105}
\problemtext 波沿$x$正方向以$u=10\,\mathrm{m/s}$传播，$\lambda=4\,\mathrm m$，$t=0$波形如图4.3.16，振幅$A=0.1\,\mathrm m$，原点$y(0,0)=0.05$ m且曲线向上，$P$点处在负半振幅$y_P=-0.05$ m的\emph{上升空间分支}。求波函数、$P$点振动方程、其坐标及回到平衡位置的最短时间。
\workings
$k=2\pi/\lambda=\pi/2\,\mathrm{rad/m}$，
$\omega=uk=5\pi\,\mathrm{rad/s}$。
令初始空间函数$y(x,0)=A\sin(kx+\varphi)$，
由$y(0,0)=A/2$且空间斜率为正，得$\varphi=\pi/6$。
\[
\boxed{y(x,t)=0.1\sin\left(\frac\pi2x-5\pi t+\frac\pi6\right)\,\mathrm m}.
\]
在$P$点，空间相位为$11\pi/6$（因$\sin=-1/2$且空间斜率为正），选离原点图示的第一个此类点：
\[
\frac\pi2x_P+\frac\pi6=\frac{11\pi}{6},
\quad
\boxed{x_P=\frac{10}{3}\,\mathrm m\approx3.333\,\mathrm m}.
\]
故
\[
\boxed{y_P(t)=0.1\sin\left(\frac{11\pi}{6}-5\pi t\right)\,\mathrm m}.
\]
$t=0$时该质点速度为负，需先经过负极值再到零。时间相位从$11\pi/6$下降至$\pi$时首次过零，变化$5\pi/6$，
\[
\boxed{t_{\min}=\frac{5\pi/6}{5\pi}=\frac16\,\mathrm s}.
\]
\crosscheck 该图形题给出的是\emph{选定的第一个上升空间分支}；若图中的$P$并非这一分支，坐标与最短时间须按实际斜率重新判断，故仍标记\textbf{待图示二校}。

\sourceproblem{4.3.17}{105}{105}
\problemtext 给定平面简谐波中$P$点振动方程
$y_P=A\cos(\omega t+\varphi_0)$；分别在原图(a)中以$O$为原点、$P$位于$x=l$且波沿$+x$传播，和图(b)中以$P$为原点、$Q$位于$x=b$且波沿$-x$传播时求波函数及$Q$点振动方程。
\workings
(a) 右行波按参考点$P$延迟时间$(x-l)/u$：
\[
\boxed{y_a(x,t)=A\cos\left[\omega t+\varphi_0-\frac{\omega}{u}(x-l)\right]}.
\]
若$Q$在$P$右边$b$处，则
$\boxed{y_{Q,a}(t)=A\cos(\omega t+\varphi_0-\omega b/u)}$。
(b) 以$P$为$x=0$，左行波使用$t+x/u$：
\[
\boxed{y_b(x,t)=A\cos\left(\omega t+\varphi_0+\frac{\omega x}{u}\right)},
\]
\[
\boxed{y_{Q,b}(t)=A\cos(\omega t+\varphi_0+\omega b/u)}.
\]
\crosscheck 对两式分别求等相位轨迹，(a)有$dx/dt=+u$、(b)有$dx/dt=-u$；代入$P$位置均恢复给定振动方程。
% === END textbook-ch04-solutions-011-017.tex ===

% === BEGIN textbook-ch04-solutions-018-023.tex blob 75ab9a3a3bb5e882560372de05f2253a763deeb4 ===
% 《大学物理简明教程》第5版，解答题4.3.18—4.3.23。
% 原书及PDF第105—106页，均依据扫描页实际图形及参数核对。
\subsection{解答题：波形追踪、干涉与相干条件}

\sourceproblem{4.3.18}{105--106}{105--106}
\problemtext 已知平面简谐波的波函数
$y=A\cos[\pi(4t+2x)]$（SI）。
(1) 写出$t=4.2\,\mathrm s$时各波峰位置、离原点最近的波峰与波峰下次经过原点的时间；
(2) 给出$t=4.2\,\mathrm s$时的波形。
\workings
$\omega=4\pi\,\mathrm{rad/s}$、$k=2\pi\,\mathrm{rad/m}$，
$\lambda=1\,\mathrm m$、$T=0.5\,\mathrm s$，
相位为$\omega t+kx$，故波沿$-x$传播，$u=2\,\mathrm{m/s}$。
波峰条件
\[
\pi(4t+2x)=2n\pi
\quad\Longrightarrow\quad
x_n(t)=n-2t,\quad n\in\mathbb Z.
\]
在$t=4.2$时
\[
\boxed{x_n=n-8.4\ \mathrm m,\quad n\in\mathbb Z}.
\]
离原点最近的波峰为$n=8$对应$x=-0.4\,\mathrm m$，
故$\boxed{|x|_{\min}=0.4\,\mathrm m}$。
波峰经过原点需$4t=2n$，即$t=n/2$；
$t>4.2$的最早值是$\boxed{t=4.5\,\mathrm s}$，即再过$0.3\,\mathrm s$；它对应$n=9$的波峰，并非已经位于原点左侧的$n=8$波峰。
固定$t=4.2$，空间波形可写为
\[
\boxed{y(x,4.2)=A\cos(0.8\pi+2\pi x)}.
\]
\crosscheck 波峰间隔1 m，随时间向左以2 m/s平移；其轨迹$dx_n/dt=-2$。

\sourceproblem{4.3.19}{106}{106}
\problemtext 如图4.3.19，波沿直线向右以$u=20\,\mathrm{m/s}$传播，四点从左到右为$C,B,A,D$，距离$CB=8$ m、$BA=5$ m、$AD=9$ m。已知$A$的振动方程$y_A=3\cos(4\pi t)$。分别以$A$、$B$为坐标原点写出波函数及$C,D$处振动方程。
\workings
$\omega=4\pi$，$k=\omega/u=\pi/5\,\mathrm{rad/m}$。
以$A$为原点、向右为正，
\[
\boxed{y_A(x,t)=3\cos\left(4\pi t-\frac\pi5x\right)}.
\]
$C$位于$x_C=-13$ m、$D$位于$x_D=9$ m，故
\[
\boxed{y_C(t)=3\cos(4\pi t+13\pi/5)},\quad
\boxed{y_D(t)=3\cos(4\pi t-9\pi/5)}.
\]
以$B$为原点，$A$位于$x'=5$ m，故
\[
\boxed{y_B(x',t)=3\cos\left(4\pi t-\frac\pi5(x'-5)\right)}
\]
（这里$y_B$表示\emph{以B为原点的整条波函数}而不专指B点）。
$C$为$x'=-8$ m、$D$为$x'=14$ m，代入重新得到上面的$y_C,y_D$。
\crosscheck 在以$A$为原点的波函数代入$x=0$恢复已给定$y_A$；坐标平移不改变任何实际质点的振动方程。振幅3的单位沿用原书未标注单位，不擅自补成3 m。

\sourceproblem{4.3.20}{106}{106}
\problemtext 两等振幅相干波源$S_1,S_2$，$S_1$至观测点$P$为$30$ cm，$S_1S_2=40$ cm且$\angle S_1=90^\circ$。在此时$S_1$位于波峰、$S_2$位于波谷，传播速率$u=100\,\mathrm{m/s}$，求在$P$发生相长干涉所需的最小频率。
\workings
由直角三角形$S_2P=50$ cm，故两路程差
$\Delta r=0.50-0.30=0.20$ m。
源在同一时刻波峰与波谷，初相差$\pi$；到达点相长条件：
\[
\pi+\frac{2\pi\Delta r}{\lambda}=2n\pi
\quad\Longrightarrow\quad
\Delta r=\left(n-\frac12\right)\lambda.
\]
为了获得\emph{最小正}频率，取满足条件的最大波长
$\lambda_{\max}=2\Delta r=0.40\,\mathrm m$，于是
\[
\boxed{\nu_{\min}=\frac{u}{\lambda_{\max}}
=\frac{100}{0.40}=250\,\mathrm{Hz}}.
\]
\crosscheck 反相波源在零程差点必相消，需额外半波长差抵消反相才能加强。

\sourceproblem{4.3.21}{106}{106}
\problemtext 两相干波源$S_1,S_2$在同一条$x$轴上，间距$d=30\,\mathrm m$；$S_1$为$x=0$，$S_2$为$x=30$ m。两波振幅相等，在$x_1=9$ m及$x_2=12$ m处为\emph{相邻}的干涉静止点，波在介质中向外传播时波长不变。求波长及两源初相差。
\workings
在源间$0<x<d$，两路程为$r_1=x,r_2=d-x$，
故路程差$\Delta r=d-2x$。沿$x$每增加$\Delta x$，干涉相位差变化$2k\Delta x$。
相邻相消位置的相位差相隔$2\pi$，因而
\[
2k(x_2-x_1)=2\pi
\quad\Longrightarrow\quad
\boxed{\lambda=2(x_2-x_1)=6\,\mathrm m}.
\]
在$x_1=9$ m，$r_2-r_1=30-18=12\,\mathrm m=2\lambda$，传播相位差为$4\pi$，模$2\pi$等于零。
为了在此处相消，两源的初相差必须为奇数倍$\pi$：
\[
\boxed{\varphi_2-\varphi_1=\pi\pmod{2\pi}}.
\]
\crosscheck 在$x_2=12$ m处路程差$6$ m$=\lambda$，同样与源反相合成相消；区间内$x=10.5$ m则由相位差额外变化$\pi$而成为相长点。

\sourceproblem{4.3.22}{106}{106}
\problemtext 图4.3.22中的等振幅相干波源$S_1,S_2$位于一条直线，间距为$\lambda/4$；$S_1$的初相位比$S_2$超前$\pi/2$，各自振幅为$A_1$。求$S_1$左侧和$S_2$右侧各点的合振幅及波强。
\workings
取由$S_1$指向$S_2$为$x$正向，$d=\lambda/4$。
对于左侧$x<0$，$r_2-r_1=d$。两波到达该点的相位差
\[
\theta_1-\theta_2
=(\varphi_1-\varphi_2)-k(r_1-r_2)
=\frac\pi2+\frac\pi2=\pi.
\]
故合振幅
$\boxed{A_{\rm left}=0}$，平均波强
$\boxed{I_{\rm left}=0}$（理想相干、等振幅）。
对于右侧$x>d$，$r_1-r_2=d$，于是
$\theta_1-\theta_2=\pi/2-\pi/2=0$，
\[
\boxed{A_{\rm right}=2A_1},\qquad
\boxed{I_{\rm right}=4I_1},
\]
其中$I_1$是单独一列波在该处传播介质中的平均波强。
\crosscheck 平均波强正比振幅平方；此结论假设在直线上两波传播到这些点的振幅相等、无额外几何衰减。

\sourceproblem{4.3.23}{106}{106}
\problemtext 两相干波源$B,C$到观测点$P$的距离分别$BP=0.40$ m、$CP=0.50$ m，传播速率$u=0.20\,\mathrm{m/s}$。
两源振动分别为$y_B=2\times10^{-3}\cos(2\pi t)$与
$y_C=2\times10^{-3}\cos(2\pi t+\pi)$（SI）。
求两波到达$P$的相位差及同时合振动的振幅。
\workings
两源频率均$1\,\mathrm{Hz}$，故波长$\lambda=u/\nu=0.20$ m，$k=2\pi/\lambda=10\pi\,\mathrm{rad/m}$。
在$P$处第二列波相对第一列波的相位差
\[
\Delta\varphi=(\pi-k CP)-(0-k BP)
=\pi-k(CP-BP)
=\pi-10\pi(0.10)=0.
\]
所以
\[
\boxed{\Delta\varphi=0\pmod{2\pi}},\qquad
\boxed{A_P=2A=4\times10^{-3}\,\mathrm m}.
\]
\crosscheck 两源虽然反相，但较长传播路径产生额外$\pi$相移，使得$P$点两波恰好同相；只有求到达点相位才能判断加强/相消。$4\times10^{-3}$ m的合振幅还要求两列波到达$P$时各自振幅均为$2\times10^{-3}$ m；若存在几何衰减或吸收，则同相时合振幅为$A_B(P)+A_C(P)$。
% === END textbook-ch04-solutions-018-023.tex ===
